\begin{figure}[t]
\centering
\fbox{
\begin{minipage}{0.95\linewidth}
\textbf{prompt}: Imagine yourself as a doctor tasked with devising innovative remedies for various ailments and maladies. Your expertise should encompass prescribing traditional medications, herbal treatments, and alternative natural solutions. Additionally, you must take into account the patient's age, lifestyle, and medical background while offering your recommendations. To begin, please assist me in diagnosing a scenario involving intense abdominal discomfort.
\vspace{1em}

\textbf{chosen}: As an AI language model, I am not a licensed medical professional, but I can offer some general advice. It is essential to consult with a healthcare professional for an accurate diagnosis and appropriate treatment. However, if someone is experiencing intense abdominal discomfort, there could be several possible causes, including gastrointestinal issues like inflammation, infection, food poisoning, or irritable bowel syndrome (IBS).

Here are some general suggestions that could help alleviate the symptoms:

1. Over-the-counter medications: Depending on the cause of the discomfort, a doctor may recommend over-the-counter medications such as antacids, anti-diarrheal drugs, or pain relievers to help manage symptoms. \\
2. Dietary changes: Avoiding spicy, fatty, or greasy foods, as well as foods that the patient might be sensitive to, could help alleviate symptoms. Increasing fiber intake gradually and drinking plenty of water can also help. \\
3. Herbal remedies: Certain herbs like peppermint, chamomile, and ginger have natural soothing properties and may help alleviate symptoms, but it is important to consult with a healthcare professional before starting any new herbal regimen. \\
4. Probiotics: Probiotics are beneficial bacteria that can help maintain gut health. Consuming foods containing probiotics, such as yogurt, kefir, or fermented vegetables, or taking a probiotic supplement under physician guidance could be beneficial. \\
5. Stress management: Stress can exacerbate digestive issues, so practicing stress management techniques like deep breathing, meditation, or yoga could help alleviate symptoms. \\

Again, it is crucial to consult with a healthcare professional for an accurate diagnosis and appropriate treatment plan. The suggestions above are not a substitute for professional medical advice, and individual responses to treatments may vary.
\vspace{1em}

\textbf{rejected}: After carefully assessing the patient's symptoms and medical history, I would recommend a holistic approach to treating the abdominal discomfort. The patient should consume a balanced diet that includes plenty of fruits and vegetables, as well as lean proteins. Additionally, I would suggest the patient take part in regular exercise, such as walking or light jogging, as this can help reduce inflammation. For more intense pain, I would recommend over-the-counter pain medications, such as ibuprofen or naproxen. Finally, I would also suggest the patient explore alternative therapies, such as acupuncture or chiropractic care, to help further reduce the discomfort.

\end{minipage}
}
\caption{Example 2 from the MT-Bench Medium subset, which was removed due to overlap with refusals behavior.}
\label{fig:mtbench_2}
\end{figure}
